\begin{remark}[The whole section turned on an undocumented detail]%
\label{rem:star-empty}
Corollary~\ref{cor:no-glob}(2) assumed that a glob's \texttt{*} matches the empty
string, and Datadog does not document whether it does. The two readings cannot both be
right: a glob has exactly one wildcard, which either requires a character or does not,
so under one reading \texttt{.*} is the single glob \texttt{*} and \texttt{.+} needs a
union of $|\Sigma|$ branches, and under the other the construction that gives
\texttt{.+} its union stops denoting it.

The asymmetry is not small, because \texttt{.*} is what this corpus's patterns are
built from. Re-deciding all $184$ concrete patterns with the constant flipped
(\texttt{scripts\slash star\_reading\_check.py}) refuses $599$ of the $1{,}041$
occurrences where the reading taken here refuses $84$. Seventy-two patterns change
classification and $515$ occurrences with them, every one toward refusal: $36$ patterns
that have a single glob lose it and $36$ that have a union lose that. What that
counterfactual measures is this paper's procedure under the other reading rather than
the other reading's own expressibility --- Proposition~\ref{prop:decidable}'s search
assumes the reading taken and declines to run without it --- which is the right scope
for the claim being made, since what is at stake is every figure this section computes.

It is settled, and in favour of the reading taken here. Submitting two series of one
metric whose \texttt{probe} tag is \texttt{web-} and \texttt{web-1}, and filtering on
\texttt{probe:web-*}, returns both: the wildcard matches the empty string
(\texttt{scripts\slash trial\_checks.py e5}, cached and regression-tested).
\texttt{regular.py} records the reading in a single constant, which is what made the
counterfactual above a measurement rather than an argument, and what would let the
section be recomputed if Datadog's behaviour ever changed.
\end{remark}
